\section{Introduction}

Symbolic regression (SR) aims to discover compact and interpretable mathematical expressions from data, making it a central tool for scientific modeling and equation discovery. While early and widely used SR systems were largely built around genetic programming and evolutionary search \citep{koza1992genetic,cranmer2023pysr}, recent methods span reinforcement learning and policy search \citep{petersen2021dsr}, modular neural-guided search \citep{landajuela2022unified}, Monte Carlo tree search and planning \citep{kamienny2023deep,shojaee2023tpsr}, transformer-based expression generation \citep{valipour2021symbolicgpt}, and LLM-assisted scientific equation discovery \citep{shojaee2025llmsr,grayeli2024lasr}. This diversity is scientifically valuable, yet it also reshapes the benchmark problem by requiring SR methods to be evaluated along multiple dimensions, including numerical accuracy, symbolic recovery, out-of-distribution generalization and seed stability.

Existing SR benchmarks provide important foundations for reproducible comparison \citep{lacava2021contemporary,matsubara2024rethinking,shojaee2025llmsrbench,xu2024geobench}. SRBench\citep{lacava2021contemporary} established a living benchmark for comparing contemporary SR methods across real-world and synthetic regression problems ; SRSD\citep{matsubara2024rethinking} revisited equation-discovery datasets with carefully designed sampling ranges, dummy variables, and expression-level similarity metrics; recent benchmarks such as LLM-SRBench\citep{shojaee2025llmsrbench} and GeoBench\citep{xu2024geobench} further target LLM-based scientific discovery and domain-specific symbolic reasoning. However, they are difficult to use as a fast, stable, and fine-grained evaluation substrate for the current algorithmic landscape. Large benchmark suites are expensive to run during iterative method development, while smaller suites can be redundant, biased toward particular formula families, or unable to preserve conclusions from a larger task pool. Static leaderboards further limit evaluation by collapsing algorithmic behavior into final aggregate scores, thereby obscuring search trajectories, invalid outputs, expression complexity, and failure modes that are central to judging practical SR utility.

We argue that a modern SR benchmark should function as evaluation infrastructure rather as a static dataset collection. It should provide reliable comparison across heterogeneous tasks, algorithms, and future submissions. %Therefore, \arena{} is designed around four principles: executable dataset standardization, compact reservoir-preserving benchmark distillation, unified algorithm execution with comparable artifacts, and an absolute dynamic leaderboard whose scores remain interpretable as new methods are added. This design enables SR evaluation to capture both final performance and algorithmic behavior, including ID/OOD accuracy, symbolic recovery, failures, complexity, and search dynamics.
%
To address these challenges, \arena{}, a standardized evaluation infrastructure for symbolic regression is proposed. It starts from nearly one thousand candidate tasks and constructs a 664-task ground-truth reservoir with standardized schemas%, executable programs, canonical expressions, ID/OOD splits, and task metadata
. A compact \core{} benchmark is distilled from the 664-task reservoir through a cascaded probe-based selection pipeline. Rather than relying on a single difficulty heuristic or manual preference, \core{} is selected by jointly optimizing task-structure coverage, response diversity, method discriminability, difficulty and failure-mode balance. 

Moreover, \arena{} introduces a fixed six-axis evaluation protocol for dynamic SR leaderboards, covering in-distribution numerical quality, out-of-distribution generalization, symbolic fidelity, efficiency cost, noise robustness, and effective stability. Combined with best-so-far search traces, this protocol supports stable long-term algorithm comparison and fine-grained analysis of SR search behavior. Empirically, we assess whether the distilled \core{} better preserves reservoir-level conclusions than common subset-selection baselines, and use the resulting leaderboard to analyze trade-offs among numerical accuracy, symbolic recovery, OOD generalization, efficiency, robustness, and stability.

Our contributions can be summarized as follows:
\begin{itemize}[leftmargin=*, itemsep=2pt, topsep=2pt, parsep=0pt]
    \item \textbf{Unified SR evaluation stack.}
    \arena{} standardizes datasets, wrappers, budgets, failure handling, expression normalization, artifacts, and minute-level logs for fair comparison across 12 SR algorithms.

    \item \textbf{Cascaded Core-50 distillation.}
    From 664 ground-truth tasks, \arena{} constructs \core{} through dual-probe mining, 12-algorithm calibration, \probe{} selection, and multi-seed reservoir probing.

    \item \textbf{Validated compact benchmark.}
    \core{} is designed to preserve coverage, diversity, discriminability, stability, difficulty balance, and low redundancy, and is validated against repeated subset-selection baselines.

    \item \textbf{Maintainable dynamic leaderboard.}
    \arena{} reports fixed-formula scores for \idq{}, \oodg{}, \symf{}, \eff{}, and \stab{}, with \rob{} as a noisy-training extension, enabling leaderboard expansion without score drift.
\end{itemize}

%The clean evaluation pipeline comprises 14,696 runs: 1,328 dual-probe runs, 2,400 \candidate{} calibration runs, 7,968 \probe{} full-reservoir runs, and 3,000 \core{} leaderboard runs. The leaderboard stage additionally produces 180,000 best-so-far expression snapshots; the optional noisy-training robustness track adds 9,000 runs. The central claim is not that 50 tasks can replace all SR evaluation. Rather, \core{} is designed as a high-information, low-cost, reproducible core benchmark that retains the main characteristics of a much larger SR task reservoir while supporting fast and extensible leaderboard updates.
